% Ad Atticum 1.1 (ad-atticum-01-01) --- English translation
% Translated by Claude (Anthropic) from the Latin (Perseus).
% Style guide: STYLE.md.

\ciceroLetterOpener{Cicero to Atticus, greetings}

\ciceroSection{1} The state of my candidacy --- which I know is your highest care --- is, so far as can be made out by conjecture, of this kind. Publius Galba alone is canvassing. He is being refused without disguise or trickery, in the old way. As men think, his over-hasty canvassing has not been against my own interest; for they refuse him in such terms as to say that they are bound to me. So I hope something is gained for us as this becomes widely known --- that we are found to have most friends. We had been thinking of beginning to canvass at the very moment when, as Cincius said, your boy was setting out with this letter, on the Field at the tribunician elections, on the seventeenth day before the Kalends of Sextilis. Competitors who appear to be standing for certain: Galba, Antonius, and Quintus Cornificius. At this you will either have laughed or groaned. To make you strike your forehead, there are some who think Caesonius too. Aquilius we did not suppose, who refused, swore he was sick, and put forward that judicial kingship of his. Catiline, if it shall have been judged that the sun does not shine at noon, will be a certain competitor. About Aufidius and Palicanus I do not suppose you are waiting for me to write.

\ciceroSection{2} Of those who are now standing, Caesar is reckoned a certainty. Thermus is thought to be contending with Silanus; who are so destitute both of friends and of standing that to me it seems not impossible \textit{[Greek: adynaton]} to bring forward Curius. But this seems so to no one but me. To our own interests it would seem most expedient that Thermus be elected with Caesar. For of those now standing, no one, if he should fall back into our year, would seem a stronger candidate, since he is curator of the Via Flaminia, which when it is finished I should be very glad to have him as my colleague\textellipsis . Of the candidates this is, so far, the rough sketch of my thinking. We shall apply the highest diligence to fulfilling every duty of a candidate, and perhaps, since Gaul seems to count for much in the voting --- when at Rome the Forum has cooled from the courts --- we shall run down in the month of September as a legate to Piso, to come back in January. When I have looked through the inclinations of the nobles, I will write to you. The rest I hope to be ample, at least so far as these urban competitors go. That other faction --- our friend Pompey's --- I want you to secure for me, since you are nearer them. Tell him that I shall not be angry with him if he does not come to my elections. So much for that.

\ciceroSection{3} But there is something for which I much wish to be pardoned by you. Your uncle Caecilius, when he was being defrauded of a great sum of money by Publius Varius, began an action against Varius's brother, Aulus Caninius Satyrus, on the ground that Satyrus had taken over property from Varius by fraud. With him the rest of the creditors were jointly suing, among whom were Lucius Lucullus and Publius Scipio and the man whom they thought would be the foreman if the goods went up for sale, Lucius Pontius. (But this point about the foreman is ridiculous.) Now learn the matter. Caecilius asked me to attend on his side against Satyrus. There is scarcely a day when this Satyrus does not come to my house: he attends most on Lucius Domitius, holds me next; he has been of great use both to me and to my brother Quintus in our candidacies.

\ciceroSection{4} I was greatly disturbed both by Satyrus's friendship and by Domitius's --- the man on whom alone our canvass most leans. I made these things plain to Caecilius, and at the same time pointed out that, had he been contending alone with Satyrus alone, I should have given him satisfaction. As things are now, in the cause of all the creditors, men of the highest standing, who without the man Caecilius would put forward in his own name could easily sustain the common case, it was fair that he should consider both my obligation and my circumstances. He took this harder, it seemed to me, than I should have wished and than well-bred men usually do, and afterwards withdrew utterly from the daily intercourse of a few days that we had begun. From you I beg pardon for this, and ask you to think that I was kept by humanity from going against the highest reputation of a friend in his bitterest hour, when he had laid all his zeal and offices upon me. If you wish to be harder on me, you will think canvassing-ambition stood in my way. But I think, even if it did, I should be pardoned --- \textit{[Greek: epei ouch hierēïon oude boeiēn]} (``since not for a sacrificial victim or an ox-hide,'' Iliad 22.159: the prize is large). For you see in what race we are running, and how we think every favour must be not only kept but even acquired. I hope I have proved my case to you; I certainly wish so.

\ciceroSection{5} Your Hermathena delights me very much, and is so prettily set up that the whole gymnasium seems an offering, \textit{[Greek: anath\=ema]}, to it. We love you greatly.
